% =========================================================
% ARTICLE — Effet de la taille de dexel sur la croissance
%            du chatter : analyse de l'amplitude exponentielle
%
% Fichier    : Etapa4_article.tex
% Auteur     : Enrique Mireles
% =========================================================

\documentclass[11pt,a4paper]{article}

% ---- Langue, encodage ----
\usepackage[utf8]{inputenc}
\usepackage[T1]{fontenc}
\usepackage{csquotes}
\usepackage[french]{babel}
\usepackage{lmodern}
\usepackage{microtype}

% ---- Mise en page ----
\usepackage[top=2.5cm,bottom=2.5cm,left=2.8cm,right=2.8cm]{geometry}
\usepackage{setspace}
\setstretch{1.12}
\setlength{\parindent}{0pt}
\setlength{\parskip}{0.55em}

% ---- Mathématiques ----
\usepackage{amsmath}
\usepackage{amssymb}
\usepackage{amsfonts}
\usepackage{mathtools}
\usepackage{bm}

% ---- Unités ----
\usepackage{siunitx}
\sisetup{output-decimal-marker={,}, separate-uncertainty=true, per-mode=symbol}

% ---- Figures ----
\usepackage{graphicx}
\usepackage{float}
\usepackage{subcaption}
\usepackage{placeins}
\graphicspath{{figures/}}

% ---- Tableaux ----
\usepackage{booktabs}
\usepackage{array}
\usepackage{multirow}

% ---- Couleurs et hyperliens ----
\usepackage{xcolor}
\definecolor{linkcolor}{RGB}{0,70,140}
\definecolor{citecolor}{RGB}{0,100,70}
\definecolor{urlcolor}{RGB}{120,40,120}

\usepackage[colorlinks=true,
            linkcolor=linkcolor,
            citecolor=citecolor,
            urlcolor=urlcolor,
            bookmarksopen=true,
            bookmarksnumbered=true]{hyperref}

\usepackage[nameinlink,noabbrev]{cleveref}
\crefname{equation}{éq.}{éqs.}
\Crefname{equation}{Éq.}{Éqs.}
\crefname{figure}{fig.}{figs.}
\Crefname{figure}{Fig.}{Figs.}
\crefname{table}{tab.}{tabs.}
\Crefname{table}{Tab.}{Tabs.}
\crefname{section}{§}{§}
\Crefname{section}{§}{§}

% ---- En-têtes ----
\usepackage{fancyhdr}
\pagestyle{fancy}
\fancyhf{}
\fancyhead[L]{\small Effet du dexel sur le chatter}
\fancyhead[R]{\thepage}
\setlength{\headheight}{13.6pt}
\renewcommand{\headrulewidth}{0.4pt}

% ---- Format des titres ----
\usepackage{titlesec}
\titleformat{\section}{\large\bfseries}{\thesection.}{0.5em}{}
\titleformat{\subsection}{\normalsize\bfseries}{\thesubsection.}{0.5em}{}
\titleformat{\subsubsection}{\normalsize\itshape}{\thesubsubsection.}{0.5em}{}

% ---- Bibliographie ----
\usepackage[backend=biber,style=numeric,sorting=none]{biblatex}
\addbibresource{references.bib}

% =========================================================
% MACROS DE NOTATION — toutes les grandeurs en un seul endroit
% =========================================================

% --- Géométrie et profondeur ---
\newcommand{\ap}{a_p}                          % profondeur de coupe
\newcommand{\apz}{a_{p,0}}                     % profondeur initiale
\newcommand{\apf}{a_{p,1}}                     % profondeur finale
\newcommand{\apcrit}{a_{p,\mathrm{crit}}}      % profondeur critique
\newcommand{\trampe}{t_{\mathrm{rampe}}}       % durée de la rampe
\newcommand{\dxl}{h_{d}}                       % taille de dexel (générique)
\newcommand{\dxli}{h_{d,i}}                    % taille de dexel du cas i

% --- Modèle dynamique ---
\newcommand{\xdyn}{x_{i,\mathrm{dyn}}}        % déplacement corrigé
\newcommand{\xdoti}{\dot{x}_{i,\mathrm{dyn}}} % vitesse corrigée
\newcommand{\omn}{\omega_n}                    % pulsation propre
\newcommand{\omnsq}{\omega_n^{2}}              % pulsation propre au carré

% --- Forces ---
\newcommand{\Fnom}{F_{\mathrm{nom}}}           % force nominale
\newcommand{\dFi}{\delta F_i}                  % force résiduelle du cas i
\newcommand{\dFin}{\delta F_{i,\mathrm{n}}}    % force résiduelle normalisée

% --- Amplitude et ajustement ---
\newcommand{\Ri}{R_i(t)}                       % RMS glissant du cas i
\newcommand{\Rik}{R_i(t_k)}                    % RMS glissant discret
\newcommand{\Tw}{T_w}                          % largeur de fenêtre
\newcommand{\betai}{\beta_i}                   % taux de croissance
\newcommand{\Pfit}{P_i^{\mathrm{fit}}}         % amplitude initiale (ajustement)
\newcommand{\cfit}{c_i^{\mathrm{fit}}}         % ln P_i^fit
% Coefficient de détermination : Q pour éviter la collision avec R_i
\newcommand{\Qi}{Q_i^{2}}

% --- Proxy de perturbation ---
\newcommand{\Pproxy}{P_{h,i}}                  % proxy d'amplitude (force)
\newcommand{\cFi}{c_i^{F}}                     % ln P_i^F
\newcommand{\RiF}{R_i^{F}(t)}                  % RMS de la force normalisée

% --- Niveau de référence et décalages ---
\newcommand{\Aref}{A_{\mathrm{ref}}}           % niveau de référence
\newcommand{\tmes}{t_i^{\mathrm{mes}}}         % temps de croisement mesuré
\newcommand{\tfit}{t_i^{\mathrm{fit}}}         % temps de croisement (ajustement)
\newcommand{\tF}{t_i^{F}}                      % temps de croisement (force)
\newcommand{\Dtmes}{\Delta t_{ij}^{\mathrm{mes}}}
\newcommand{\Dtfit}{\Delta t_{ij}^{\mathrm{fit}}}
\newcommand{\DtF}{\Delta t_{ij}^{F}}
\newcommand{\DtPfit}{\Delta t_{P,ij}^{\mathrm{fit}}}
\newcommand{\DtPF}{\Delta t_{P,ij}^{F}}
\newcommand{\Dtbfit}{\Delta t_{\beta,ij}^{\mathrm{fit}}}
\newcommand{\DtbF}{\Delta t_{\beta,ij}^{F}}

% Opérateur RMS
\newcommand{\RMS}{\operatorname{RMS}}

% =========================================================
% COMMANDE POUR LES BLOCS DE FIGURES RÉSERVÉES
% =========================================================

\newcommand{\figplaceholder}[2]{%
    \begin{figure}[H]
        \centering
        \fbox{\parbox[c][#1][c]{0.88\textwidth}{%
            \centering\small #2}}
    \end{figure}%
}

% =========================================================
% DOCUMENT
% =========================================================

\title{%
    \textbf{Effet de la taille de dexel sur la croissance du chatter :}\\
    \textbf{analyse de l'amplitude exponentielle et décomposition}\\
    \textbf{du décalage temporel}%
}
\author{Enrique Mireles}
\date{\today}

\begin{document}

\maketitle

% =========================================================
\begin{abstract}
Dans les simulations de fraisage sur cône tronqué, une discrétisation plus
fine de la géométrie (dexel plus petit) retarde systématiquement le développement
exponentiel du chatter. Ce travail présente une méthode quantitative pour
décomposer ce retard en deux contributions distinctes : une différence du niveau
initial de perturbation, liée à l'amplitude de la réponse en début de zone
instable, et une différence du taux de croissance exponentielle, liée au
comportement du système dans la zone instable. Pour chaque taille de dexel,
l'évolution de l'amplitude RMS glissante est ajustée par un modèle log-linéaire,
permettant d'extraire un taux de croissance \(\betai\) et une amplitude initiale
\(\Pfit\). La force résiduelle normalisée est utilisée comme proxy observable de
la perturbation initiale associée à la discrétisation, offrant une estimation
indépendante de l'amplitude. Trois définitions complémentaires du décalage
temporel sont établies et comparées, puis décomposées analytiquement sous deux
hypothèses : taux de croissance identiques, et cas général de taux différents.
\end{abstract}

\tableofcontents
\newpage

% =========================================================
\section{Introduction}
\label{sec:introduction}
% =========================================================

Dans les simulations de fraisage, la géométrie de la pièce est représentée par
une discrétisation spatiale en \emph{dexels} de taille \(\dxl\). Cette résolution
conditionne la précision avec laquelle la surface usinée — et donc l'épaisseur
instantanée de matière enlevée — est représentée entre deux passages successifs
de l'outil.

L'observation centrale de cette étude est la suivante : lorsque \(\dxl\)
diminue, le début de la croissance exponentielle du déplacement dynamique est
décalé vers des temps plus élevés. En désignant par \(t_{\mathrm{exp},i}\)
l'instant à partir duquel le cas \(i\) présente un comportement exponentiel, on
observe

\begin{equation}
    \dxli < h_{d,j}
    \quad\Longrightarrow\quad
    t_{\mathrm{exp},i} > t_{\mathrm{exp},j},
    \label{eq:observation}
\end{equation}

c'est-à-dire qu'un dexel plus fin retarde l'apparition de la croissance
exponentielle. Cette relation décrit une observation numérique ; elle n'identifie
pas encore le mécanisme responsable.

Le décalage entre deux courbes peut provenir de deux mécanismes distincts,
non exclusifs :

\begin{enumerate}
    \item \textbf{Différence du niveau initial de perturbation.} La
    discrétisation introduit une erreur géométrique qui génère une perturbation
    lors du premier contact. Si cette perturbation est plus faible pour un dexel
    fin, la réponse exponentielle démarre à un niveau plus bas et atteint un
    niveau donné plus tard.
    \item \textbf{Différence du taux de croissance.} La représentation de la
    surface influe sur l'épaisseur instantanée de coupe et donc sur le couplage
    régénératif. Si ce couplage est modifié par \(\dxl\), le taux de croissance
    \(\betai\) peut différer entre les cas.
\end{enumerate}

L'objectif de cet article est de séparer et de quantifier ces deux contributions.
La méthode repose sur l'ajustement log-linéaire de l'amplitude RMS glissante,
la définition d'un niveau de référence commun, et l'utilisation de la force
résiduelle normalisée comme proxy observable de la perturbation initiale due à
la discrétisation.

La \ref{sec:setup} décrit le modèle numérique et les paramètres de simulation.
La \ref{sec:pretreat} présente les corrections appliquées aux signaux bruts.
La \ref{sec:amplitude} définit la grandeur d'amplitude et les trois
estimateurs du décalage temporel. Les \ref{sec:h1} et \ref{sec:h2} analysent
respectivement l'hypothèse de taux identiques et le cas général.
La \ref{sec:beta-origin} discute l'origine physique des différences de taux de
croissance.


% =========================================================
\section{Configuration numérique}
\label{sec:setup}
% =========================================================

\subsection{Géométrie et profondeur de coupe variable}
\label{subsec:geometrie}

La pièce est un cône tronqué usiné à profondeur de coupe croissante. Dans le
repère temporel de la simulation, \(\ap(t)\) évolue linéairement entre une
valeur initiale \(\apz\) et une valeur finale \(\apf\) :

\begin{equation}
    \ap(t)
    =
    \apz + \left(\apf-\apz\right)\frac{t}{\trampe},
    \qquad 0 \leq t \leq \trampe,
    \label{eq:rampe}
\end{equation}

avec \(\apz = \SI{5}{\milli\metre}\) et \(\apf = \SI{15}{\milli\metre}\).

\begin{figure}[H]
    \centering
    \fbox{\parbox[c][4cm][c]{0.88\textwidth}{\centering\small
        \textbf{[Fig. 1 — à ajouter : schéma manuel]}\\[0.2cm]
        Cône tronqué avec la direction d'usinage, $\apz$, $\apf$
        et la longueur de la zone conique.
    }}
    \caption{Géométrie de la pièce et rampe de profondeur de coupe.}
    \label{fig:cone}
\end{figure}

\subsection{Modèle dynamique régénératif}
\label{subsec:model}

La dynamique de l'outil dans la direction axiale est représentée par un
oscillateur modal soumis à une force régénérative :

\begin{equation}
    m\ddot{x}(t) + c\dot{x}(t) + kx(t)
    =
    K_f\,\ap(t)\bigl[x(t-T)-x(t)\bigr],
    \label{eq:edm}
\end{equation}

où \(m\), \(c\), \(k\) sont la masse, l'amortissement et la rigidité modaux,
\(K_f\) le coefficient de force, \(T\) la période de révolution, et
\(x(t-T)-x(t)\) le terme régénératif. La profondeur de coupe \(\ap(t)\) étant
variable, l'équation est à coefficients non constants. Le système évolue d'une
zone stable (\(\ap < \apcrit\)) vers une zone instable (\(\ap > \apcrit\)) au
cours de la simulation, avec

\begin{equation}
    \apz < \apcrit < \apf.
    \label{eq:traversee}
\end{equation}

L'épaisseur instantanée de matière enlevée est

\begin{equation}
    h(t) = f_z + x(t-T) - x(t),
    \label{eq:thickness}
\end{equation}

où \(f_z\) est l'avance nominale. La force de l'équation~\eqref{eq:edm}
correspond donc à \(K_f\ap(t)[x(t-T)-x(t)]\), c'est-à-dire la seule
contribution dynamique (régénérative) de la force.

\begin{figure}[H]
    \centering
    \fbox{\parbox[c][3.5cm][c]{0.88\textwidth}{\centering\small
        \textbf{[Fig. 2 — à générer depuis le code]}\\[0.2cm]
        Diagramme des lóbulos de stabilité avec la vitesse de rotation étudiée,
        $\apcrit$ et la rampe de $\apz$ à $\apf$.
    }}
    \caption{Position de la rampe dans le diagramme de stabilité.}
    \label{fig:lobes}
\end{figure}

\subsection{Cas simulés et tailles de dexel}
\label{subsec:dexels}

Chaque simulation utilise une taille de dexel \(\dxli\) différente. Une
valeur plus faible de \(\dxli\) correspond à une représentation géométrique
plus fine de la surface usinée. Les \(N\) cas étudiés sont notés

\begin{equation}
    i \in \{1, \ldots, N\},
    \qquad h_{d,1} > h_{d,2} > \cdots > h_{d,N}.
    \label{eq:ordering-dexel}
\end{equation}

Pour chaque cas \(i\), les signaux extraits de la simulation sont :

\begin{equation}
    x_i(t),
    \qquad \dot{x}_i(t),
    \qquad F_i(t),
    \label{eq:signals}
\end{equation}

où \(x_i(t)\) est le déplacement axial, \(\dot{x}_i(t)\) la vitesse axiale et
\(F_i(t)\) la force calculée dans la direction étudiée.

\subsection{Paramètres numériques}
\label{subsec:params}

Les principaux paramètres sont regroupés dans le \ref{tab:params}.

\begin{table}[H]
    \centering
    \caption{Paramètres du modèle numérique.}
    \label{tab:params}
    \begin{tabular}{llll}
        \toprule
        Paramètre & Symbole & Valeur & Unité \\
        \midrule
        Vitesse de rotation        & \(n\)         & \(12094{,}28\)          & tr/min \\
        Avance nominale            & \(f_z\)        & \(0{,}05\)              & \si{\milli\metre} \\
        Coefficient de force       & \(K_f\)        & \(1{,}0\times10^{9}\)   & \si{\newton\per\metre\squared} \\
        Profondeur initiale        & \(\apz\)       & \(5\)                   & \si{\milli\metre} \\
        Profondeur finale          & \(\apf\)       & \(15\)                  & \si{\milli\metre} \\
        Longueur de la zone conique & \(L\)         & \(150\)                 & \si{\milli\metre} \\
        Rigidité modale            & \(k\)          & \(2{,}13\times10^{8}\)  & \si{\newton\per\metre} \\
        Fréquence propre           & \(f_n\)        & \(150\)                 & \si{\hertz} \\
        Pulsation propre           & \(\omn\)       & \(942{,}48\)            & \si{\radian\per\second} \\
        Pulsation propre au carré  & \(\omnsq = k/m\)   & \(8{,}88\times10^{5}\) & \si{\per\second\squared} \\
        Masse modale               & \(m = k/\omnsq\)   & \(239{,}79\)           & \si{\kilogram} \\
        Angle de projection        & \(\theta\)    & \(135\)                 & \si{\degree} \\
        Coeff.\ de projection modale & \(\phi=\sin\theta/\sqrt{m}\) & \(4{,}57\times10^{-2}\) & \si{\per\sqrt{\kilogram}} \\
        \bottomrule
    \end{tabular}
\end{table}

\begin{figure}[H]
    \centering
    \fbox{\parbox[c][4cm][c]{0.88\textwidth}{\centering\small
        \textbf{[Fig. 3 — à générer depuis le code]}\\[0.2cm]
        Panneaux \(x_i(t)\), \(\dot{x}_i(t)\), \(F_i(t)\) pour chaque taille
        de dexel (arrangement \(3\times N\) ou \(N\) colonnes de 3 lignes).\\
        But : montrer le décalage visuel de la croissance exponentielle.
    }}
    \caption{Signaux simulés pour les différentes tailles de dexel.}
    \label{fig:signals}
\end{figure}


% =========================================================
\section{Prétraitement des signaux}
\label{sec:pretreat}
% =========================================================

La profondeur de coupe croissante \(\ap(t)\) introduit une tendance
déterministe dans la force, le déplacement et la vitesse. Ces tendances sont
soustraites pour isoler la réponse dynamique pure, qui seule porte l'information
sur la croissance du chatter.

\subsection{Force nominale et résiduelle}
\label{subsec:fnominal}

La force attendue en l'absence d'oscillations est

\begin{equation}
    \Fnom(t) = K_f\,\ap(t)\,f_z.
    \label{eq:Fnom}
\end{equation}

La force résiduelle du cas \(i\) est la déviation par rapport à cette valeur :

\begin{equation}
    \dFi(t) = F_i(t) - \Fnom(t).
    \label{eq:Fresid}
\end{equation}

La quantité \(\dFi(t)\) contient les variations de force qui ne sont pas
expliquées par la profondeur de coupe nominale. D'après le modèle régénératif,
elle s'écrit

\begin{equation}
    \dFi(t) \approx K_f\,\ap(t)\bigl[x_i(t-T) - x_i(t)\bigr].
    \label{eq:Fresid-regen}
\end{equation}

\begin{figure}[H]
    \centering
    \fbox{\parbox[c][3.8cm][c]{0.88\textwidth}{\centering\small
        \textbf{[Fig. 4 — à générer depuis le code]}\\[0.2cm]
        Pour chaque cas \(i\) : \(F_i(t)\), \(\Fnom(t)\) et \(\dFi(t)\)
        superposés. But : vérifier la décomposition et montrer le niveau
        du résidu.
    }}
    \caption{Décomposition de la force en composante nominale et résiduelle.}
    \label{fig:Fresid}
\end{figure}

\subsection{Correction du déplacement}
\label{subsec:xcorr}

La force nominale produit une déflexion quasi statique de l'outil. Dans la
formulation modale, cette déflexion est

\begin{equation}
    \delta_{\mathrm{stat}}(t)
    =
    \frac{\phi^{2}}{\omnsq}\,\Fnom(t)
    =
    \frac{\phi^{2}}{\omnsq}\,K_f\,\ap(t)\,f_z.
    \label{eq:delta-stat}
\end{equation}

Comme \(\ap(t)\) est linéaire (cf.\ \ref{eq:rampe}), \(\delta_{\mathrm{stat}}\)
est également linéaire. Le déplacement dynamique corrigé est

\begin{equation}
    \xdyn(t) = x_i(t) - \delta_{\mathrm{stat}}(t).
    \label{eq:xdyn}
\end{equation}

\subsection{Correction de la vitesse}
\label{subsec:vdotcorr}

La dérivée de \(\delta_{\mathrm{stat}}\) est constante car \(\ap(t)\) est
linéaire :

\begin{equation}
    \dot{\delta}_{\mathrm{stat}}
    =
    \frac{\phi^{2}}{\omnsq}\,K_f\,f_z\,
    \frac{\apf - \apz}{\trampe}.
    \label{eq:ddot-stat}
\end{equation}

La vitesse corrigée est

\begin{equation}
    \xdoti(t) = \dot{x}_i(t) - \dot{\delta}_{\mathrm{stat}}.
    \label{eq:xdotdyn}
\end{equation}

\begin{figure}[H]
    \centering
    \fbox{\parbox[c][3.8cm][c]{0.88\textwidth}{\centering\small
        \textbf{[Fig. 5 — à générer depuis le code]}\\[0.2cm]
        Pour un cas représentatif : \(x_i(t)\) (brut), \(\delta_{\mathrm{stat}}(t)\)
        et \(\xdyn(t)\) (corrigé). Idem pour la vitesse.
    }}
    \caption{Correction de la composante quasi statique du déplacement et de la vitesse.}
    \label{fig:xcorr}
\end{figure}

\subsection{Normalisation de la force résiduelle}
\label{subsec:fnorm}

D'après \ref{eq:Fresid-regen}, le terme régénératif $[x_i(t-T)-x_i(t)]$
est multiplié par \(\ap(t)\). Deux perturbations évaluées à des instants
différents de la rampe ont donc des amplitudes différentes uniquement à cause
de l'évolution de \(\ap(t)\), et non de la discrétisation. Pour supprimer cet
effet direct, on normalise :

\begin{equation}
    \dFin(t) = \frac{\dFi(t)}{\ap(t)}.
    \label{eq:Fnorm}
\end{equation}

En substituant \ref{eq:Fresid-regen},

\begin{equation}
    \dFin(t) \approx K_f\bigl[x_i(t-T) - x_i(t)\bigr].
    \label{eq:Fnorm-regen}
\end{equation}

La dépendance directe à \(\ap(t)\) est ainsi supprimée. La grandeur \(\dFin(t)\)
est proportionnelle à la différence régénérative de déplacement, et sera
utilisée en \ref{subsec:proxy} comme proxy observable de la perturbation
initiale associée à la taille de dexel.

\begin{figure}[H]
    \centering
    \fbox{\parbox[c][3.8cm][c]{0.88\textwidth}{\centering\small
        \textbf{[Fig. 6 — à générer depuis le code]}\\[0.2cm]
        Comparaison de \(\dFi(t)\) et \(\dFin(t)\) pour les différents cas.
        But : montrer que la normalisation supprime la tendance croissante.
    }}
    \caption{Effet de la normalisation de la force résiduelle.}
    \label{fig:Fnorm}
\end{figure}


% =========================================================
\section{Analyse de l'amplitude exponentielle}
\label{sec:amplitude}
% =========================================================

\subsection{Amplitude RMS glissante}
\label{subsec:rms}

Pour chaque signal corrigé \(z_i(t)\) (déplacement \(\xdyn\), vitesse \(\xdoti\),
ou force normalisée \(\dFin\)), l'amplitude instantanée est estimée par la
valeur efficace sur une fenêtre centrée glissante de durée \(\Tw\) :

\begin{equation}
    R_i(t) =
    \sqrt{
        \frac{1}{\Tw}
        \int_{t-\Tw/2}^{t+\Tw/2} z_i^2(\tau)\,\mathrm{d}\tau
    }.
    \label{eq:rms-cont}
\end{equation}

En pratique, la discrétisation donne

\begin{equation}
    \Rik =
    \sqrt{
        \frac{1}{N_w}
        \sum_{n\in\mathcal{W}_k} z_i^2(t_n)
    },
    \label{eq:rms-disc}
\end{equation}

où \(N_w\) est le nombre d'échantillons dans la fenêtre \(\mathcal{W}_k\). La
largeur de fenêtre est choisie pour couvrir plusieurs périodes de la réponse
dominante, sans supprimer l'évolution temporelle de l'enveloppe :

\begin{equation}
    \Tw = \frac{5}{f_n}.
    \label{eq:Tw}
\end{equation}

\subsection{Modèle log-linéaire et ajustement}
\label{subsec:logfit}

Dans la zone instable, l'amplitude suit une croissance exponentielle. En
désignant par \(I_i = [t_{a,i}, t_{b,i}]\) l'intervalle où cette hypothèse
est vérifiée pour le cas \(i\), le modèle est

\begin{equation}
    R_i(t) = \Pfit\,\exp(\betai t), \qquad t \in I_i.
    \label{eq:exp-model}
\end{equation}

En prenant le logarithme naturel,

\begin{equation}
    \ln R_i(t) = \cfit + \betai t, \qquad t \in I_i,
    \label{eq:loglinear}
\end{equation}

avec \(\cfit = \ln\Pfit\). La pente \(\betai\) (\si{\per\second}) est le taux de
croissance exponentielle et \(\Pfit\) est l'amplitude extrapolée à \(t=0\) à
partir de l'ajustement. Ces deux paramètres sont estimés par régression linéaire
de \(\ln R_i(t)\) sur \(I_i\).

La qualité de l'ajustement est évaluée par le coefficient de détermination

\begin{equation}
    \Qi =
    1 -
    \frac{
        \sum_{k}\bigl(\ln R_i(t_k) - \cfit - \betai t_k\bigr)^{2}
    }{
        \sum_{k}\bigl(\ln R_i(t_k) - \overline{\ln R_i}\bigr)^{2}
    },
    \label{eq:Q2}
\end{equation}

où \(\overline{\ln R_i}\) est la moyenne sur \(I_i\). Un \(\Qi \approx 1\) valide
l'hypothèse exponentielle dans l'intervalle choisi.

\begin{figure}[H]
    \centering
    \fbox{\parbox[c][4cm][c]{0.88\textwidth}{\centering\small
        \textbf{[Fig. 7 — à générer depuis le code]}\\[0.2cm]
        \(\ln R_i(t)\) pour tous les cas, avec les intervalles \(I_i\) surlignés,
        les droites ajustées, les valeurs \(\betai\) et \(\Qi\).\\
        But : montrer que les pentes sont proches ou différentes selon les cas.
    }}
    \caption{Ajustement log-linéaire des amplitudes RMS pour les différentes tailles de dexel.}
    \label{fig:logfit}
\end{figure}

\subsection{Niveau de référence commun}
\label{subsec:Aref}

Pour comparer les différents cas à un même instant d'amplitude, on choisit un
niveau de référence \(\Aref\) atteint par toutes les courbes dans leurs
intervalles d'ajustement respectifs :

\begin{equation}
    \ln\Aref
    \;\in\;
    \bigcap_{i=1}^{N}
    \bigl[
        \min_{t\in I_i}\ln R_i(t),\;
        \max_{t\in I_i}\ln R_i(t)
    \bigr].
    \label{eq:Aref-condition}
\end{equation}

\(\Aref\) n'est pas un seuil physique ; il sert uniquement à définir une
référence commune pour les mesures de décalage.

\subsection{Définition des trois estimateurs de décalage}
\label{subsec:three-dt}

Trois estimateurs indépendants du décalage temporel entre les cas \(i\) et
\(j\) sont introduits. Ils seront comparés pour vérifier la cohérence de
l'analyse.

\subsubsection{Décalage mesuré directement}

Le temps \(\tmes\) auquel la courbe RMS du cas \(i\) atteint \(\Aref\) est

\begin{equation}
    R_i\!\left(\tmes\right) = \Aref.
    \label{eq:tmes-def}
\end{equation}

Le décalage mesuré entre les cas \(i\) et \(j\) est

\begin{equation}
    \Dtmes = t_j^{\mathrm{mes}} - \tmes.
    \label{eq:Dtmes}
\end{equation}

Convention : \(\Dtmes > 0\) signifie que le cas \(j\) atteint \(\Aref\) après
le cas \(i\), c'est-à-dire que \(j\) est retardé par rapport à \(i\).

\subsubsection{Décalage issu de l'ajustement}

À partir du modèle \ref{eq:loglinear}, le temps d'atteinte de \(\Aref\) est

\begin{equation}
    \tfit = \frac{\ln\Aref - \cfit}{\betai},
    \label{eq:tfit}
\end{equation}

et le décalage prédit par les ajustements est

\begin{equation}
    \Dtfit = t_j^{\mathrm{fit}} - \tfit
    =
    \frac{\ln\Aref - c_j^{\mathrm{fit}}}{\beta_j}
    -
    \frac{\ln\Aref - \cfit}{\betai}.
    \label{eq:Dtfit-general}
\end{equation}

\subsubsection{Décalage estimé par le proxy de force}
\label{subsec:proxy}

La force résiduelle normalisée \(\dFin(t)\) est utilisée comme
\textbf{proxy observable de la perturbation initiale} associée à la
discrétisation. L'idée est la suivante : la taille de dexel détermine l'erreur
géométrique sur la surface usinée, qui se traduit directement dans le terme
régénératif de la force. Cette perturbation ne peut pas être mesurée
directement (elle est liée à la différence entre la surface réelle et sa
représentation discrète), mais elle est accessible à travers \(\dFin(t)\), qui
lui est proportionnelle d'après \ref{eq:Fnorm-regen}.

Le proxy d'amplitude \(\Pproxy\) est la valeur RMS locale de \(\dFin(t)\) à
l'instant \(t_{s,i}\) situé dans l'intervalle d'ajustement :

\begin{equation}
    \Pproxy = \RMS\!\left[\dFin(t)\right]_{t=t_{s,i}}.
    \label{eq:proxy-def}
\end{equation}

En substituant \ref{eq:Fnorm-regen},

\begin{equation}
    \Pproxy \approx K_f\,\RMS\!\left[x_i(t-T)-x_i(t)\right]_{t=t_{s,i}}.
    \label{eq:proxy-regen}
\end{equation}

Dans l'intervalle d'ajustement, le déplacement \(x_i(t)\) est dominé par une
réponse oscillatoire d'enveloppe exponentielle. Les deux termes \(x_i(t-T)\) et
\(x_i(t)\) partagent le même taux de croissance \(\betai\), de sorte que leur
différence possède également cette enveloppe. On peut donc écrire

\begin{equation}
    \RMS\!\left[x_i(t-T)-x_i(t)\right] \approx C_i\,R_i^F(t),
    \label{eq:proxy-Ci}
\end{equation}

où \(C_i\) est un facteur de proportionnalité dépendant principalement du retard
\(T\), de la fréquence dominante et de la relation de phase entre \(x_i(t-T)\)
et \(x_i(t)\), et où \(R_i^F(t)\) désigne l'amplitude RMS de \(\dFin(t)\). Le proxy
devient donc

\begin{equation}
    \Pproxy \approx K_f\,C_i\,R_i^F(t_{s,i}).
    \label{eq:proxy-approx}
\end{equation}

\paragraph{Expression analytique de \(C_i\).}
Pour un signal de la forme \(x_i(t)=A_i e^{\betai t}\sin(\omega_i t+\varphi_i)\),
la différence régénérative vaut

\[
    x_i(t-T)-x_i(t)
    =A_i e^{\betai t}\!\left[e^{-\betai T}\sin(\omega_i t-\omega_i T+\varphi_i)
    -\sin(\omega_i t+\varphi_i)\right].
\]

En calculant son RMS sur plusieurs cycles, on obtient l'expression exacte

\begin{equation}
    C_i = \sqrt{1 + e^{-2\betai T} - 2\,e^{-\betai T}\cos(\omega_i T)}.
    \label{eq:Ci-exact}
\end{equation}

Comme \(\betai T\ll 1\) (taux de croissance de quelques \si{\per\second},
retard \(T\) de l'ordre de la milliseconde pour \(N_z=1\) dent),
\(e^{-\betai T}\approx 1\) et l'expression se réduit à

\begin{equation}
    C_i \approx 2\left|\sin\!\left(\frac{\omega_i T}{2}\right)\right|.
    \label{eq:Ci-approx}
\end{equation}

La fréquence dominante du chatter est proche de la fréquence propre du système,
ici \(\omega_i \approx 2\pi\times\SI{150}{\hertz}\), identique pour tous les
cas. Avec le même retard \(T\), on a donc \(C_i\approx C_j\), ce qui légitime
l'hypothèse utilisée dans les sections suivantes.

\paragraph{Correction temporelle du proxy.}
L'amplitude extrapolée à \(t=0\) issue de l'ajustement est \(\Pfit\), tandis que
\(\Pproxy\) est évalué à l'instant \(t_{s,i}\). Pour les rendre comparables, on
applique la correction temporelle du modèle exponentiel :

\begin{equation}
    R_i^F(t_{s,i}) = P_i^F\,\exp(\betai\,t_{s,i}),
    \label{eq:RiF-ts}
\end{equation}

où \(P_i^F\) est l'amplitude extrapolée à \(t=0\) à partir de la force. En
substituant dans \ref{eq:proxy-approx} et en isolant \(P_i^F\),

\begin{equation}
    P_i^F
    =
    \frac{\Pproxy}{K_f\,C_i}\,\exp(-\betai\,t_{s,i}).
    \label{eq:PiF}
\end{equation}

En posant \(\cFi = \ln P_i^F\), il vient

\begin{equation}
    \cFi
    =
    \ln\Pproxy - \ln K_f - \ln C_i - \betai\,t_{s,i}.
    \label{eq:cF}
\end{equation}

Le temps d'atteinte de \(\Aref\) reconstruit à partir de la force est alors

\begin{equation}
    \tF
    =
    \frac{\ln\Aref - \cFi}{\betai}
    =
    t_{s,i} + \frac{1}{\betai}\ln\!\left(\frac{K_f\Aref}{\Pproxy}\right),
    \label{eq:tF}
\end{equation}

et le décalage correspondant est

\begin{equation}
    \DtF = t_j^F - \tF.
    \label{eq:DtF-general}
\end{equation}

\begin{figure}[H]
    \centering
    \fbox{\parbox[c][4cm][c]{0.88\textwidth}{\centering\small
        \textbf{[Fig. 8 — à générer depuis le code]}\\[0.2cm]
        \(\ln R_i(t)\) pour tous les cas avec la ligne horizontale \(\ln\Aref\),
        les temps \(\tmes\), \(\tfit\) et \(\tF\) marqués pour chaque cas.\\
        But : montrer la cohérence des trois estimateurs.
    }}
    \caption{Niveau de référence $\Aref$ et les trois temps d'atteinte par cas.}
    \label{fig:Aref}
\end{figure}

\begin{figure}[H]
    \centering
    \fbox{\parbox[c][3.8cm][c]{0.88\textwidth}{\centering\small
        \textbf{[Fig. 9 — à générer depuis le code]}\\[0.2cm]
        \(R_i^F(t)\) (RMS de \(\dFin(t)\)) pour chaque cas, avec \(t_{s,i}\)
        et \(\Pproxy\) marqués. But : visualiser le proxy et son instant d'évaluation.
    }}
    \caption{Proxy de perturbation issu de la force résiduelle normalisée.}
    \label{fig:proxy}
\end{figure}

\begin{figure}[H]
    \centering
    \fbox{\parbox[c][3.8cm][c]{0.88\textwidth}{\centering\small
        \textbf{[Fig. 10 — à générer depuis le code]}\\[0.2cm]
        Graphique de barres : \(\Dtmes\), \(\Dtfit\), \(\DtF\) pour chaque
        paire \((i,j)\). But : vérifier que les trois estimateurs sont cohérents.
    }}
    \caption{Comparaison des trois estimateurs du décalage temporel.}
    \label{fig:dt-compare}
\end{figure}


% =========================================================
\section{Hypothèse~1 : taux de croissance identiques}
\label{sec:h1}
% =========================================================

La première hypothèse suppose que la taille de dexel modifie le niveau initial
de perturbation mais pas le taux de croissance. Pour deux cas \(i\) et \(j\),
cela s'écrit

\begin{equation}
    \betai = \beta_j = \beta.
    \label{eq:h1}
\end{equation}

Sous cette hypothèse, les droites ajustées dans le plan \(\ln R\) vs \(t\)
sont parallèles et le décalage horizontal entre elles est constant, quelle que
soit la valeur de \(\Aref\).

\subsection{Simplification de \texorpdfstring{$\Dtfit$}{Delta t fit}}

En imposant \ref{eq:h1} dans \ref{eq:Dtfit-general},

\begin{align}
    \Dtfit
    &=
    \frac{\ln\Aref - c_j^{\mathrm{fit}}}{\beta}
    -
    \frac{\ln\Aref - \cfit}{\beta}
    \nonumber\\
    &=
    \frac{\cfit - c_j^{\mathrm{fit}}}{\beta}
    =
    \frac{1}{\beta}\ln\!\left(\frac{\Pfit}{P_j^{\mathrm{fit}}}\right).
    \label{eq:Dtfit-h1}
\end{align}

Le niveau \(\Aref\) disparaît : le décalage ne dépend que du rapport des
amplitudes initiales extrapolées.

\subsection{Simplification de \texorpdfstring{$\DtF$}{Delta t F}}

En substituant \ref{eq:tF} avec \(\betai = \beta_j = \beta\), et en regroupant
les logarithmes,

\begin{align}
    \DtF
    &=
    (t_{s,j} - t_{s,i})
    +
    \frac{1}{\beta}\ln\!\left(\frac{K_f\Aref}{P_{h,j}}\right)
    -
    \frac{1}{\beta}\ln\!\left(\frac{K_f\Aref}{\Pproxy}\right)
    \nonumber\\
    &=
    (t_{s,j} - t_{s,i})
    +
    \frac{1}{\beta}\ln\!\left(\frac{\Pproxy}{P_{h,j}}\right).
    \label{eq:DtF-h1-general}
\end{align}

Les termes \(K_f\) et \(\Aref\) s'annulent. Si de plus les facteurs de
proportionnalité \(C_i \approx C_j\), la dépendance aux \(C_i\) disparaît
également :

\begin{equation}
    \DtF
    \approx
    (t_{s,j} - t_{s,i})
    +
    \frac{1}{\beta}\ln\!\left(\frac{\Pproxy}{P_{h,j}}\right).
    \label{eq:DtF-h1}
\end{equation}

Le premier terme corrige la différence entre les instants d'évaluation des
proxies ; le second représente le décalage dû au rapport des niveaux de
perturbation mesurés à partir de la force.

\subsection{Évaluation de l'hypothèse}
\label{subsec:eval-h1}

L'hypothèse~1 est validée si :
\begin{enumerate}
    \item les \(\betai\) extraits de l'ajustement sont effectivement proches
    entre eux ;
    \item \(\Dtmes \approx \Dtfit\), c'est-à-dire que le décalage observé est
    bien reproduit par les amplitudes initiales extrapolées avec une pente
    commune ;
    \item \(\Dtfit \approx \DtF\), c'est-à-dire que la différence entre les
    amplitudes initiales est cohérente avec les niveaux de perturbation estimés
    à partir de la force.
\end{enumerate}

Si ces conditions ne sont pas toutes vérifiées, il faut abandonner l'hypothèse~1
et traiter le cas général.

\begin{figure}[H]
    \centering
    \fbox{\parbox[c][4cm][c]{0.88\textwidth}{\centering\small
        \textbf{[Fig. 11 — à générer depuis le code]}\\[0.2cm]
        Graphique de barres : \(\Dtmes\), \(\Dtfit\) (éq.~\eqref{eq:Dtfit-h1})
        et \(\DtF\) (éq.~\eqref{eq:DtF-h1}) sous l'hypothèse~1.\\
        But : valider ou invalider l'hypothèse de pentes communes.
    }}
    \caption{Validation de l'hypothèse de taux de croissance identiques.}
    \label{fig:h1-eval}
\end{figure}


% =========================================================
\section{Cas général : taux de croissance différents}
\label{sec:h2}
% =========================================================

Lorsque les pentes \(\betai\) et \(\beta_j\) obtenues par ajustement présentent
une différence non négligeable, le décalage temporel dépend à la fois du niveau
initial et du taux de croissance. On abandonne \ref{eq:h1} ; aucune contrainte
n'est imposée sur les \(\betai\).

\subsection{Décomposition analytique de \texorpdfstring{$\Dtfit$}{Delta t fit}}
\label{subsec:decomp-fit}

L'expression générale \ref{eq:Dtfit-general} est réécrite en ajoutant et
soustrayant le terme \((\ln\Aref - \cfit)/\beta_j\) :

\begin{align}
    \Dtfit
    &=
    \underbrace{
        \frac{\cfit - c_j^{\mathrm{fit}}}{\beta_j}
    }_{\DtPfit}
    +
    \underbrace{
        \left(\ln\Aref - \cfit\right)
        \left(\frac{1}{\beta_j} - \frac{1}{\betai}\right)
    }_{\Dtbfit}.
    \label{eq:Dtfit-decomp}
\end{align}

\paragraph{Contribution des amplitudes initiales.}

\begin{equation}
    \DtPfit
    =
    \frac{1}{\beta_j}
    \ln\!\left(\frac{\Pfit}{P_j^{\mathrm{fit}}}\right).
    \label{eq:DtP-fit}
\end{equation}

C'est le décalage qui serait observé si les deux ajustements avaient le même
taux de croissance \(\beta_j\) mais des niveaux initiaux différents.

\paragraph{Contribution des taux de croissance.}

\begin{equation}
    \Dtbfit
    =
    \ln\!\left(\frac{\Aref}{\Pfit}\right)
    \left(\frac{1}{\beta_j} - \frac{1}{\betai}\right).
    \label{eq:Dtb-fit}
\end{equation}

Le premier facteur, \(\ln(\Aref/\Pfit)\), est la distance logarithmique entre
le niveau de référence et l'amplitude initiale du cas \(i\). Le second facteur
mesure la différence entre les inverses des taux de croissance. Cette
contribution est nulle si et seulement si \(\betai = \beta_j\), ce qui redonne
le cas de l'hypothèse~1.

\paragraph{Cohérence de la décomposition.} Par construction,

\begin{equation}
    \Dtfit = \DtPfit + \Dtbfit.
    \label{eq:decomp-check-fit}
\end{equation}

Cette égalité constitue un critère de vérification numérique.

\subsection{Décomposition analytique de \texorpdfstring{$\DtF$}{Delta t F}}
\label{subsec:decomp-force}

En utilisant \(\cFi\) (défini en \ref{eq:cF}) comme ordonnée à l'origine
reconstruite depuis la force, le décalage \(\DtF\) a la même structure
que \(\Dtfit\). En ajoutant et soustrayant \((\ln\Aref - \cFi)/\beta_j\),
on obtient

\begin{align}
    \DtF
    &=
    \underbrace{
        \frac{\cFi - c_j^{F}}{\beta_j}
    }_{\DtPF}
    +
    \underbrace{
        \left(\ln\Aref - \cFi\right)
        \left(\frac{1}{\beta_j} - \frac{1}{\betai}\right)
    }_{\DtbF}.
    \label{eq:DtF-decomp}
\end{align}

En développant \(\cFi\) et \(c_j^{F}\) à partir de \ref{eq:cF},

\begin{align}
    \DtPF
    &=
    \frac{1}{\beta_j}
    \ln\!\left(
        \frac{\Pproxy\cdot C_j}{P_{h,j}\cdot C_i}
    \right)
    -
    \frac{\betai}{\beta_j}\,t_{s,i}
    +
    t_{s,j}.
    \label{eq:DtP-force}
\end{align}

Cette expression contient trois effets : (1) le rapport des proxies
\(\Pproxy/P_{h,j}\), (2) la différence éventuelle entre les facteurs \(C_i\) et
\(C_j\), et (3) la correction des instants d'évaluation \(t_{s,i}\),
\(t_{s,j}\).

La contribution des taux de croissance est

\begin{align}
    \DtbF
    &=
    \Bigl(
        \ln\Aref
        - \ln\Pproxy
        + \ln K_f
        + \ln C_i
        + \betai\,t_{s,i}
    \Bigr)
    \left(\frac{1}{\beta_j} - \frac{1}{\betai}\right).
    \label{eq:Dtb-force}
\end{align}

Si \(C_i \approx C_j\), le rapport de ces facteurs dans \ref{eq:DtP-force}
vaut~1 et la contribution d'amplitude se simplifie en

\begin{equation}
    \DtPF
    \approx
    \frac{1}{\beta_j}
    \ln\!\left(\frac{\Pproxy}{P_{h,j}}\right)
    -
    \frac{\betai}{\beta_j}\,t_{s,i}
    +
    t_{s,j}.
    \label{eq:DtP-force-simple}
\end{equation}

\paragraph{Cohérence de la décomposition.}

\begin{equation}
    \DtF = \DtPF + \DtbF.
    \label{eq:decomp-check-force}
\end{equation}

\subsection{Interprétation}
\label{subsec:interp-h2}

La décomposition sépare deux causes du décalage :

\begin{description}
    \item[Effet du niveau initial (\(\DtPfit\), \(\DtPF\))]
    associé à la différence entre les amplitudes extrapolées à \(t=0\). Un
    dexel plus fin génère une perturbation régénérative initiale plus faible, ce
    qui retarde l'atteinte de \(\Aref\).

    \item[Effet du taux de croissance (\(\Dtbfit\), \(\DtbF\))]
    associé à la différence entre les \(\betai\). Sa magnitude dépend à la fois
    de \(|\beta_j - \betai|\) et de la distance logarithmique
    \(\ln(\Aref/\Pfit)\). Elle peut être dominante si les taux diffèrent
    significativement.
\end{description}

La comparaison entre les contributions issues des ajustements et celles issues
de la force permet de vérifier si la différence de niveau initial est bien
expliquée par la perturbation géométrique de la discrétisation.

\begin{figure}[H]
    \centering
    \fbox{\parbox[c][4cm][c]{0.88\textwidth}{\centering\small
        \textbf{[Fig. 12 — à générer depuis le code]}\\[0.2cm]
        Graphique de barres empilées pour chaque paire \((i,j)\) :
        \(\DtPfit + \Dtbfit = \Dtfit\) vs \(\Dtmes\), et idem pour la force.\\
        But : montrer quelle contribution domine et vérifier la cohérence.
    }}
    \caption{Décomposition du décalage en contributions de niveau initial et de taux de croissance.}
    \label{fig:decomp}
\end{figure}


% =========================================================
\section{Origine des différences de taux de croissance}
\label{sec:beta-origin}
% =========================================================

\textbf{[Section à compléter — en attente du code et de la démonstration.]}

Les sections précédentes ont établi une méthode pour mesurer et décomposer le
décalage temporel entre cas de dexel différents, sans expliquer pourquoi les
taux de croissance \(\betai\) peuvent différer. Cette section a pour objectif de
relier \(\betai\) aux paramètres de la discrétisation.

Les points suivants seront développés :

\begin{enumerate}
    \item \textbf{Lien entre \(\betai\) et \(\ap(t)\) local.} Dans la
    formulation régénérative linéarisée, le taux de croissance instantané
    dépend de la profondeur de coupe effective. Si la discrétisation introduit
    une erreur sur \(\ap\), elle peut modifier \(\betai\).

    \item \textbf{Erreur géométrique introduite par le dexel.} L'erreur entre
    la surface réelle et sa représentation discrète génère une profondeur de
    coupe effective différente de la profondeur nominale. Cette erreur est bornée
    par \(\dxl\) et dépend de la géométrie du cône.

    \item \textbf{Relation entre erreur géométrique et \(\Delta\beta\).}
    La différence \(\beta_i - \beta_j\) est reliée à la différence entre les
    profondeurs de coupe effectives des deux cas dans la zone instable.

    \item \textbf{Vérification numérique.} Les \(\betai\) extraits par
    ajustement seront comparés aux valeurs prédites par le modèle analytique.
\end{enumerate}


% =========================================================
\section{Conclusion}
\label{sec:conclusion}
% =========================================================

Cet article a présenté une méthode systématique pour analyser et décomposer le
retard temporel introduit par le raffinement de la discrétisation dexel dans
les simulations de chatter. Les contributions principales sont les suivantes.

\begin{enumerate}
    \item \textbf{Trois estimateurs complémentaires} du décalage temporel ont
    été définis : le décalage mesuré directement (\(\Dtmes\)), le décalage
    prédit par les paramètres d'ajustement (\(\Dtfit\)), et le décalage
    reconstruit à partir de la force résiduelle normalisée (\(\DtF\)). Leur
    cohérence constitue un premier critère de validation.

    \item \textbf{La force résiduelle normalisée} est utilisée comme proxy
    observable de la perturbation initiale due à la discrétisation. Cette
    perturbation géométrique ne peut pas être mesurée directement, mais la
    force \(\dFin(t)\) lui est proportionnelle d'après le modèle régénératif.

    \item \textbf{Sous l'hypothèse de taux identiques} (\(\betai = \beta_j\)),
    le décalage ne dépend que du rapport des amplitudes initiales. Le niveau
    de référence \(\Aref\) disparaît de l'expression, ce qui rend la mesure
    robuste au choix de ce niveau.

    \item \textbf{Dans le cas général}, le décalage est décomposé
    analytiquement en une contribution de niveau initial (\(\DtPfit\)) et une
    contribution de taux de croissance (\(\Dtbfit\)). Cette décomposition
    permet d'identifier le mécanisme dominant pour chaque paire de cas.

    \item \textbf{L'origine physique} des différences de \(\betai\) fera
    l'objet d'une analyse complémentaire (\ref{sec:beta-origin}), reliant
    l'erreur géométrique introduite par le dexel à la modification du taux de
    croissance dans la zone instable.
\end{enumerate}

% \printbibliography

\end{document}
